/chapter{Monero 加密货币 (Monero the Cryptocurrency)}
/label{chapter:Monero-cryptocurrency}

最终章节将所有内容串联在一起（？技术性足够强到可以跳过报告其余部分？可能使篇幅过长——目标是引导读者查阅报告的其他部分）


-曲线 ed25519
-地址创建
    -基本地址
    -子地址
    -多重签名
    -支付 ID（集成地址）
-接收旧交易
    -扫描区块链，查看 output 地址，使用 view key 和交易公钥查看是否有匹配的 spend key
    -整理结果
-构建交易
    -交易结构（数据的序列化/组织和传输方式）
    -获取预期接收者的地址（或子地址），指定每个接收者的金额
    -根据发送金额 + 手续费，选择一组 owned output 使得 sum(a)>=(sum(b)+fee)，并设置找零 = sum(a) - (sum(b)+fee)
    -构建 output：
        -创建交易公钥（如果有至少一个子地址且 output 数 >1 则创建多个）
        -为每个 output 创建一次性 output 地址
        -为每个 output 金额创建承诺 C(b)
        -为每个 output 承诺创建 Diffie-Hellman 'amount' 和 'mask' 项
        -范围证明：对每个 output 金额进行 borromean 环签名
        -如有支付 ID，将其加密
    -构建 input：
        -如果交易类型为 rcttypesimple 则创建伪 output 承诺
        -为每个 input 的 MLSAG 选择环成员
        -为每个 input 计算 key image
        -在每个 input 上构建 MLSAG 签名，签名的消息包含所有其他交易信息
    -适当地填写交易数据结构（贯穿交易全过程）
-提交交易
    -验证后放入内存池
-挖矿入区块链
    -交易被组织成 Merkle 树并进行哈希
    -搜索 nonce 直到满足难度要求
    -区块发布到网络
    -区块被接受或拒绝（被孤立）；共识机制
    -[可能的] 签名数据被修剪
